\chapter{Functional Requirements Catalogue}
\label{app:requirements}

This appendix is the full functional requirements catalogue summarised in
Section~\ref{sec:fr}. Requirements are atomic and testable. Each states input, precondition,
behaviour, output, and error or rejection behaviour, and carries a priority
(Must, Should, Could). The identifier prefix names the capability:
\texttt{FR-CLI} (command line and run control), \texttt{FR-UI} (web user
interface), \texttt{FR-VER} (verification engine), \texttt{FR-SBX} (sandbox),
\texttt{FR-AGENT} (agents and model provider), \texttt{FR-EVO} (evolutionary
loop), \texttt{FR-PER} (persistence and recovery), and \texttt{FR-HAR}
(harness). The acceptance-evidence column names real tests (from the committed
test suite) or campaign artifacts; the testing chapter maps these to the
\texttt{TC-*} families.

\subsection{Command line and run control}
\label{sec:fr-cli}

\begin{longtable}{@{}p{0.10\textwidth}p{0.40\textwidth}p{0.07\textwidth}p{0.33\textwidth}@{}}
\caption{Functional requirements: command line and run control.}\label{tab:fr-cli}\\
\toprule
\textbf{ID} & \textbf{Requirement} & \textbf{Priority} & \textbf{Acceptance evidence} \\
\midrule
\endfirsthead
\toprule
\textbf{ID} & \textbf{Requirement} & \textbf{Priority} & \textbf{Acceptance evidence} \\
\midrule
\endhead
\bottomrule
\endfoot

\texttt{FR-CLI-01} & \textbf{Configure and validate a run.} Input: a
configuration (CLI flags or YAML). Precondition: none. Behaviour: build a
\texttt{RunConfig}, validate every field (problem, provider, population 1 to
50, generations 1 to 1000, timeout 1 to 300 s, bound sizes, budgets), and
resolve a problem pack when one is named. Output: a saved
\texttt{config.yaml} under the run directory. Error: on any invalid or unknown
field, list every error and exit non-zero; do not start a run. & Must &
\texttt{test\_config\_pack\_and\_problem\_round\_trip},
\texttt{test\_config\_with\_bad\_pack\_reports\_error},
\texttt{test\_config\_validation\_for\_api}, \texttt{test\_unknown\_key\_rejected} \\
\addlinespace[2pt]

\texttt{FR-CLI-02} & \textbf{Launch a run.} Input: a validated configuration.
Precondition: \texttt{FR-CLI-01} succeeded. Behaviour: initialize the
population, then run the evolutionary loop until a stop condition, then grade
the best passing candidates on the final bound. Output: per-generation progress
and artifacts. Error: a single failed model call is recorded and the loop
continues; an invalid configuration aborts before any run starts. & Must &
\texttt{test\_launch}, \texttt{test\_launch\_from\_pack},
\texttt{test\_launch\_with\_api\_endpoint}, \texttt{test\_mock\_run\_end\_to\_end} \\
\addlinespace[2pt]

\texttt{FR-CLI-03} & \textbf{Report run status.} Input: a run directory.
Precondition: the run exists. Behaviour: read the latest checkpoint and the
engine liveness marker and report generation, best fitness, best work ratio,
pass rate, and recent failures. Output: a status summary (text or JSON). Error:
an unknown run is reported and the command exits non-zero. & Must &
\texttt{test\_list\_and\_get\_run},
\texttt{test\_engine\_liveness\_and\_interrupted\_state} \\
\addlinespace[2pt]

\texttt{FR-CLI-04} & \textbf{Export results.} Input: a run directory.
Precondition: at least one generation completed. Behaviour: write a
machine-readable \texttt{run\_summary.json} and an
\texttt{efficiency\_report.csv} with per-candidate fitness, work ratio, state,
and test counts. Output: the two files under the run's results directory. Error:
with no results, report it and exit non-zero. & Must &
\texttt{test\_final\_and\_finish}; campaign \texttt{run\_summary.json} and
\texttt{efficiency\_report.csv} \\
\addlinespace[2pt]

\texttt{FR-CLI-05} & \textbf{Review run history.} Input: a run directory root.
Precondition: at least one prior run exists. Behaviour: list runs with their
summaries (problem, generations, best fitness, outcome). Output: a history
listing (text or JSON). Error: with no runs, state that and exit zero. & Should &
\texttt{test\_list\_and\_get\_run} \\
\addlinespace[2pt]

\texttt{FR-CLI-06} & \textbf{Steer a running run.} Input: a control command
(\texttt{pause}, \texttt{resume}, \texttt{stop}, \texttt{hint}, \texttt{set})
and its arguments. Precondition: the run is active. Behaviour: append the
command to the run's control file; the orchestrator applies it at the next
control point and acknowledges it. Output: a \texttt{control\_ack} or
\texttt{control\_reject} event. Error: a command for an ended run is refused; an
unknown command is rejected at validation; each command is applied exactly once
across restarts. & Must & \texttt{test\_steering\_hint\_and\_population},
\texttt{test\_pause\_applies\_before\_next\_agent\_call},
\texttt{test\_pause\_blocks\_until\_resume},
\texttt{test\_commands\_are\_handled\_once\_across\_restarts},
\texttt{test\_control\_refused\_on\_ended\_run},
\texttt{test\_send\_command\_rejects\_unknown} \\
\addlinespace[2pt]

\texttt{FR-CLI-07} & \textbf{Resume a run.} Input: a run directory. Precondition:
at least one committed checkpoint exists. Behaviour: reload the latest complete
checkpoint, restore the specialist state from the event log, reseed parent
selection for the current generation, and continue. Output: the run resumes from
the saved generation without redoing completed work. Error: a corrupt or
malformed checkpoint is refused and reported; the last good generation is used. &
Must & \texttt{test\_resume\_keeps\_call\_count},
\texttt{test\_resume\_draws\_the\_same\_selection\_stream},
\texttt{test\_resume\_counts\_calls\_of\_a\_crashed\_generation},
\texttt{test\_malformed\_state\_fails\_cleanly} \\
\addlinespace[2pt]

\texttt{FR-CLI-08} & \textbf{Run the verifier self-test.} Input: a problem and a
graph-size bound. Precondition: the verifier is built. Behaviour: run the
known-good references, the seeded fault corpus, and the equivalent controls
through the full verification path. Output: a pass/fail verdict with the
zero-false-negative count. Error: if a known-good is rejected or a fault is
accepted, the self-test fails and no experiment may rely on the verifier. & Must &
\texttt{test\_known\_good\_references\_pass\_n5},
\texttt{test\_real\_faults\_rejected\_n5},
\texttt{test\_equivalent\_controls\_accepted\_n5},
\texttt{test\_yang\_dfs\_reference\_passes\_n5}; campaign
\texttt{selftest\_dfs\_n5.log}, \texttt{selftest\_lbfs\_n5.log} \\
\addlinespace[2pt]

\texttt{FR-CLI-09} & \textbf{Serve the web UI.} Input: a run directory root, a
bind host, and a port. Precondition: none. Behaviour: start the local web server
bound to the loopback address by default, generate an access token, and print a
token link once. Output: a running server serving the UI and the run API. Error:
every API request and websocket without the token is refused. & Must &
\texttt{test\_access\_token}, \texttt{test\_cross\_site\_origin\_refused} \\
\bottomrule
\end{longtable}

\subsection{Web user interface}
\label{sec:fr-ui}

\begin{longtable}{@{}p{0.10\textwidth}p{0.40\textwidth}p{0.07\textwidth}p{0.33\textwidth}@{}}
\caption{Functional requirements: web user interface.}\label{tab:fr-ui}\\
\toprule
\textbf{ID} & \textbf{Requirement} & \textbf{Priority} & \textbf{Acceptance evidence} \\
\midrule
\endfirsthead
\toprule
\textbf{ID} & \textbf{Requirement} & \textbf{Priority} & \textbf{Acceptance evidence} \\
\midrule
\endhead
\bottomrule
\endfoot

\texttt{FR-UI-01} & \textbf{Configure and launch from the UI.} Input: the New
Run form (problem pack, generations, population, call budget, provider and
model, verifier bound, run id, dynamic-agents toggle). Precondition: the server
is running. Behaviour: validate the form client-side and launch a run through
the API. Output: a new run that appears in the run list. Error: out-of-range
numbers, bad run ids, a missing pack, or an API endpoint without a model and key
are blocked with named messages. & Must & vitest \texttt{newrun.test.ts}:
``a default form can launch'', ``blocks out-of-range numbers, bad run ids and a
missing pack'', ``warns when the call budget cannot cover the planned calls'',
``an API endpoint needs a model and a key''; \texttt{ui\_check/01\_runs.png} \\
\addlinespace[2pt]

\texttt{FR-UI-02} & \textbf{Monitor a run live.} Input: a run. Precondition: the
run is active. Behaviour: subscribe to the event stream and render the current
generation, best fitness and work ratio, the population, and recent failures.
Output: a live run view. Error: an offline or interrupted run is shown with its
state, not as a crash. & Must & vitest \texttt{state.test.ts}: ``builds
generations, candidates and status''; \texttt{test\_run\_emits\_full\_event\_stream};
\texttt{ui\_check/04\_live.png} \\
\addlinespace[2pt]

\texttt{FR-UI-03} & \textbf{Steer a run from the UI.} Input: a control command
from the Steer panel (hint, set population or budget, pause, resume, stop).
Precondition: the run is active. Behaviour: send the command over the websocket;
show the acknowledgement. Output: the engine applies the command at the next
control point. Error: a command for an ended run is shown as refused. & Must &
\texttt{test\_steering\_hint\_and\_population}; vitest \texttt{status.test.ts}:
``pause: queued to pausing to paused'', ``stop and reject clear pending'';
\texttt{ui\_check/05\_steer.png} \\
\addlinespace[2pt]

\texttt{FR-UI-04} & \textbf{Replay generation history.} Input: a finished run.
Precondition: at least two generations exist. Behaviour: a generation scrubber
steps through committed generations, updating the population and metrics for the
selected generation. Output: the replay reflects each generation's state. Error:
the scrubber stops at the last complete generation. & Must &
\texttt{test\_websocket\_backlog\_from\_offset\_and\_live};
\texttt{ui\_check/07\_replay\_g2.png} \\
\addlinespace[2pt]

\texttt{FR-UI-05} & \textbf{Show failure patterns.} Input: a run with failures.
Precondition: at least one candidate failed. Behaviour: count the failed
candidates of each generation by kind (wrong after init, insert, delete, crash,
timeout) and by the size of the first failing graph; each bar segment opens an
example candidate. Output: a navigable failure-patterns panel. Error: with no
failures, show an empty state. & Must & vitest \texttt{state.test.ts}: ``counts
failures per generation by kind and graph size''; \texttt{ui\_check/08\_failures.png},
\texttt{ui\_check/09\_failure\_drawer.png} \\
\addlinespace[2pt]

\texttt{FR-UI-06} & \textbf{Show efficiency and lineage charts.} Input: a run
with metrics. Precondition: at least two generations with work ratios. Behaviour:
render fitness-over-generation, work-ratio-over-generation, and candidate-lineage
charts. Output: rendered charts with accessible labels. Error: if a metric is
absent, omit that chart. & Should & \texttt{ui\_check/06\_done.png}
(\texttt{charts.tsx} \texttt{aria-label}) \\
\addlinespace[2pt]

\texttt{FR-UI-07} & \textbf{Inspect a candidate.} Input: a candidate id.
Precondition: the candidate exists. Behaviour: open a drawer showing the
candidate's code, its state, fitness, work ratio, and its failing case. Output:
the candidate detail. Error: an unknown or path-traversal candidate id returns
404. & Must & \texttt{test\_unknown\_or\_traversal\_is\_404};
\texttt{ui\_check/09\_failure\_drawer.png} \\
\addlinespace[2pt]

\texttt{FR-UI-08} & \textbf{Authenticate UI access.} Input: the access token.
Precondition: the server is running. Behaviour: accept the token via an
\texttt{HttpOnly} cookie (set from the printed link) or a bearer header; refuse
requests carrying a non-local browser Origin. Output: authorized access or a
403. Error: a missing or wrong token is refused with a message pointing at the
printed link. & Must & \texttt{test\_access\_token},
\texttt{test\_cross\_site\_origin\_refused},
\texttt{test\_dashboard\_escapes\_script\_end\_in\_llm\_text} \\
\addlinespace[2pt]

\texttt{FR-UI-09} & \textbf{Provide an accessible interface.} Input: none.
Precondition: none. Behaviour: use semantic controls with ARIA roles and labels
(\texttt{role="dialog"} for the candidate drawer, \texttt{role="status"}
\texttt{aria-live} for activity, \texttt{aria-current} for the run stepper,
\texttt{aria-label} for charts), and move keyboard focus to the relevant field on
step change. Output: a keyboard-navigable, screen-reader-labelled UI. Error: full
WCAG conformance requires manual testing with assistive technology and is not
asserted by the automated suite. & Should & \texttt{ui/src/components/*}
(\texttt{aria-*}, \texttt{role=*}, focus management); vitest \texttt{state.test.ts} \\
\bottomrule
\end{longtable}

\subsection{Verification engine}
\label{sec:fr-ver}

\begin{longtable}{@{}p{0.10\textwidth}p{0.40\textwidth}p{0.07\textwidth}p{0.33\textwidth}@{}}
\caption{Functional requirements: verification engine.}\label{tab:fr-ver}\\
\toprule
\textbf{ID} & \textbf{Requirement} & \textbf{Priority} & \textbf{Acceptance evidence} \\
\midrule
\endfirsthead
\toprule
\textbf{ID} & \textbf{Requirement} & \textbf{Priority} & \textbf{Acceptance evidence} \\
\midrule
\endhead
\bottomrule
\endfoot

\texttt{FR-VER-01} & \textbf{Isomorphism-aware enumeration.} Input: a graph-size
bound and a problem. Precondition: the \texttt{nauty} data files are present.
Behaviour: for directed problems, load the non-isomorphic digraph
representatives (rooted classes for rooted problems) from the committed data
files and add deterministic relabelings of each representative; for undirected
problems, use every labelled graph up to the bound. Output: a cached set of base
graphs per size. Error: the universe is deterministic across calls; a missing
data file falls back to seeded random samples for that size. & Must &
\texttt{test\_class\_counts\_directed}, \texttt{test\_class\_counts\_rooted},
\texttt{test\_universe\_deterministic\_across\_calls},
\texttt{test\_relabel\_determinism} \\
\addlinespace[2pt]

\texttt{FR-VER-02} & \textbf{Exhaustive single-edge flips.} Input: a base graph
of size at most the exhaustive-flip bound (5). Precondition: none. Behaviour:
for every legal edge, emit a length-1 sequence that inserts it if absent or
deletes it if present, so every possible single update is tested. Output: the
flip sequences for the graph. Error: an operation on a loop or an already
illegal edge is never emitted. & Must &
\texttt{test\_fault\_counterexamples\_trigger\_invariant},
\texttt{test\_universe\_deterministic\_across\_calls} \\
\addlinespace[2pt]

\texttt{FR-VER-03} & \textbf{Seeded random sequences.} Input: a base graph, a
sequence count, a maximum length, and a seed. Precondition: none. Behaviour:
generate the configured number of legal operation sequences of up to the maximum
length, seeded by the bound seed plus the graph size, so the universe is
reproducible. Output: the random sequences. Error: every operation stays legal
on the evolving graph. & Must & \texttt{test\_universe\_relabels\_scale},
\texttt{test\_relabel\_cap\_at\_n2},
\texttt{test\_universe\_deterministic\_across\_calls} \\
\addlinespace[2pt]

\texttt{FR-VER-04} & \textbf{Invariant checking after every operation.} Input: a
candidate state after each operation and the current graph. Precondition: the
problem defines an invariant or certificate. Behaviour: check the invariant after
the initial state and after every operation; record the first violation with its
graph, operation, and context. Output: pass/fail per sequence with the violating
operation. Error: a malformed candidate state fails that candidate without
aborting the run. & Must & \texttt{test\_bfs\_distance\_field\_is\_certified},
\texttt{test\_dfs\_rejects\_parent\_that\_did\_not\_discover},
\texttt{test\_dfs\_requires\_finish},
\texttt{test\_dfs\_checker\_rejects\_nan\_parent\_cycle\_without\_hanging},
\texttt{test\_scc\_rejects\_split\_cycle\_through\_three\_components},
\texttt{test\_none\_distances\_pass} \\
\addlinespace[2pt]

\texttt{FR-VER-05} & \textbf{Fitness scoring.} Input: per-sequence pass/fail
results. Precondition: the candidate was executed. Behaviour: set fitness to the
fraction of sequences whose invariant held throughout. Output: a fitness value in
0 to 1. Error: if no sequence could be executed, fitness is 0. & Must &
\texttt{test\_best\_is\_best\_pass\_so\_far} \\
\addlinespace[2pt]

\texttt{FR-VER-06} & \textbf{Efficiency scoring.} Input: candidate work and
offline-recomputation work. Precondition: the candidate ran to completion.
Behaviour: compute the work ratio as candidate work divided by offline work, where
work counts the instrumented update steps; an adaptive repeat keeps the timed call
total under a budget. Output: a work-ratio value, or none when the candidate did
not complete. Error: a candidate that did not complete has no work ratio. & Must &
\texttt{test\_signature\_without\_failure\_records\_is\_still\_correctness};
campaign \texttt{analysis.md} (work-ratio definition) \\
\addlinespace[2pt]

\texttt{FR-VER-07} & \textbf{Efficiency-improvement flag.} Input: a work ratio
and a fitness. Precondition: the candidate ran. Behaviour: mark the candidate an
efficiency improvement only if it is fully correct (fitness 1) and its work ratio
is below 1. Output: the improvement flag. Error: a work ratio of 1 or more means
no improvement; it is not an error. & Must &
\texttt{test\_best\_is\_best\_pass\_so\_far} (ranking key uses work ratio) \\
\addlinespace[2pt]

\texttt{FR-VER-08} & \textbf{Verifier self-test.} Input: known-good references,
a seeded fault corpus, and equivalent controls. Precondition: the verifier is
built. Behaviour: accept every known-good reference and equivalent control, and
reject every fault, on the stated bound. Output: a verdict with the
zero-false-negative count. Error: any accepted fault or rejected reference fails
the self-test. & Must & \texttt{test\_known\_good\_references\_pass\_n5},
\texttt{test\_real\_faults\_rejected\_n5},
\texttt{test\_equivalent\_controls\_accepted\_n5},
\texttt{test\_fault\_corpus\_metadata\_complete},
\texttt{test\_yang\_dfs\_reference\_and\_mutations\_in\_process}; campaign
\texttt{selftest\_dfs\_n5.log}, \texttt{selftest\_lbfs\_n5.log} \\
\addlinespace[2pt]

\texttt{FR-VER-09} & \textbf{Trust boundary.} Input: none. Precondition: none.
Behaviour: the verifier, the test graphs, and the oracle are never readable or
writable by a candidate or an agent; agents receive only fitness and structured
failure summaries. Output: enforced isolation. Error: N/A (enforced by
construction and the sandbox). & Must & \texttt{test\_candidate\_cannot\_reach\_secrets\_or\_network},
\texttt{test\_harness\_env\_is\_allow\_listed} \\
\addlinespace[2pt]

\texttt{FR-VER-10} & \textbf{Two-tier grading.} Input: the population and the
bound. Precondition: a search completed. Behaviour: grade every candidate on the
small search bound during the run, then grade the best passing candidates on the
full final bound, trying the next-best if one fails. Output: search results kept
and a final-grade verdict. Error: a search pass that fails the final bound is
reported, and the next passing candidate is tried. & Must &
\texttt{test\_final\_grade\_tries\_next\_passing},
\texttt{test\_final\_and\_finish}; campaign \texttt{analysis.md} eval1 \\
\addlinespace[2pt]

\texttt{FR-VER-11} & \textbf{Update-mode filtering.} Input: an update mode
(\texttt{mixed}, \texttt{insert\_only}, \texttt{delete\_only}). Precondition:
none. Behaviour: for a restricted mode, replay each sequence on a copy of the base
graph and keep only the legal operations of the allowed type, so semi-dynamic
problems can be declared without changing the generators. Output: legal sequences
for the mode. Error: an invalid mode raises at construction; a restricted sequence
stays legal. & Should & \texttt{test\_update\_mode\_insert\_only},
\texttt{test\_update\_mode\_delete\_only},
\texttt{test\_update\_mode\_mixed\_unchanged},
\texttt{test\_update\_mode\_sequences\_stay\_legal},
\texttt{test\_update\_mode\_invalid\_raises} \\
\bottomrule
\end{longtable}

\subsection{Sandbox}
\label{sec:fr-sbx}

\begin{longtable}{@{}p{0.10\textwidth}p{0.40\textwidth}p{0.07\textwidth}p{0.33\textwidth}@{}}
\caption{Functional requirements: sandbox.}\label{tab:fr-sbx}\\
\toprule
\textbf{ID} & \textbf{Requirement} & \textbf{Priority} & \textbf{Acceptance evidence} \\
\midrule
\endfirsthead
\toprule
\textbf{ID} & \textbf{Requirement} & \textbf{Priority} & \textbf{Acceptance evidence} \\
\midrule
\endhead
\bottomrule
\endfoot

\texttt{FR-SBX-01} & \textbf{Namespace jail.} Input: a candidate source and the
interpreter. Precondition: Linux with user namespaces. Behaviour: run the
candidate in a new user, network (down), mount, pid, and IPC namespace over a
read-only file tree that exposes only the interpreter and its import paths, with
host data directories hidden and a fresh \texttt{/proc} and \texttt{/tmp}.
Output: an isolated execution environment. Error: the candidate source travels in
the input, so the run directory is never visible inside. & Must &
\texttt{test\_jail\_fails\_closed},
\texttt{test\_unjailed\_forked\_child\_cannot\_hold\_the\_run\_open} \\
\addlinespace[2pt]

\texttt{FR-SBX-02} & \textbf{Resource limits.} Input: a per-chunk timeout.
Precondition: none. Behaviour: before candidate code runs, set the address-space
limit to 4 GB, the CPU limit to the timeout plus 5 s, the process count to 4 in
the runner's user namespace, the per-file limit to 16 MB, the open-file limit to
64, and the result output to 256 MB. Output: a candidate that cannot exhaust the
host. Error: a candidate that raises a limit is stopped, not allowed. & Must &
\texttt{test\_candidate\_cannot\_exhaust\_host\_resources},
\texttt{test\_limits\_leave\_room\_for\_a\_final\_grade\_chunk},
\texttt{test\_cpu\_limit\_stops\_a\_spinning\_candidate},
\texttt{test\_shared\_memory\_does\_not\_outlive\_the\_candidate},
\texttt{test\_sandbox\_leaves\_no\_temp\_files} \\
\addlinespace[2pt]

\texttt{FR-SBX-03} & \textbf{Timeout and hang location.} Input: a candidate and a
per-chunk wall-clock timeout. Precondition: none. Behaviour: kill a candidate that
exceeds the timeout and record it as a failure; after a timeout, locate the first
hanging graph and sequence with a bounded probe so the reflector gets a concrete
case. Output: a timeout failure with a located hang. Error: a timeout is a
failure, never a pass, and never halts the run. & Must &
\texttt{test\_infinite\_loop\_times\_out\_serial},
\texttt{test\_infinite\_loop\_times\_out\_parallel}, \texttt{test\_hang\_is\_located},
\texttt{test\_killed\_candidate\_is\_reported\_as\_kill\_and\_traceback\_has\_line} \\
\addlinespace[2pt]

\texttt{FR-SBX-04} & \textbf{No secrets or network for candidates.} Input: a
candidate. Precondition: none. Behaviour: scrub the environment, hide the
verifier's credentials and host files, and keep the network namespace down, so a
candidate cannot read secrets or reach the network. Output: a candidate with no
access to credentials or the network. Error: a candidate that attempts either is
blocked and the run continues. & Must &
\texttt{test\_candidate\_cannot\_reach\_secrets\_or\_network},
\texttt{test\_agent\_options\_never\_leak\_keys},
\texttt{test\_api\_error\_text\_drops\_response\_body} \\
\addlinespace[2pt]

\texttt{FR-SBX-05} & \textbf{Fail closed.} Input: the host capabilities.
Precondition: none. Behaviour: if the jail cannot be built (not Linux, user
namespaces disabled), refuse to run candidate code unless explicitly disabled.
Output: no unconfined execution by default. Error: disabling the jail is explicit
and not the default. & Must & \texttt{test\_jail\_fails\_closed} \\
\bottomrule
\end{longtable}

\subsection{Agents and model provider}
\label{sec:fr-agent}

\begin{longtable}{@{}p{0.10\textwidth}p{0.40\textwidth}p{0.07\textwidth}p{0.33\textwidth}@{}}
\caption{Functional requirements: agents and model provider.}\label{tab:fr-agent}\\
\toprule
\textbf{ID} & \textbf{Requirement} & \textbf{Priority} & \textbf{Acceptance evidence} \\
\midrule
\endfirsthead
\toprule
\textbf{ID} & \textbf{Requirement} & \textbf{Priority} & \textbf{Acceptance evidence} \\
\midrule
\endhead
\bottomrule
\endfoot

\texttt{FR-AGENT-01} & \textbf{Provider abstraction.} Input: a provider
(\texttt{api}, \texttt{enlcode}, or \texttt{mock}) and a model. Precondition: the
provider is configured. Behaviour: route every model call through a single
vendor-agnostic backend interface; no component depends on a specific vendor.
Output: model responses. Error: an unknown provider is rejected at configuration
time; a missing API key is a clear error. & Must &
\texttt{test\_mock\_run\_end\_to\_end}, \texttt{test\_launch\_with\_api\_endpoint},
\texttt{test\_enlcode\_command\_and\_parse},
\texttt{test\_api\_missing\_key\_is\_a\_clear\_error} \\
\addlinespace[2pt]

\texttt{FR-AGENT-02} & \textbf{Generation role.} Input: a problem definition, a
pack brief, and an optional parent. Precondition: the provider is reachable.
Behaviour: produce a new candidate implementing the candidate interface
(\texttt{init}, \texttt{insert\_edge}, \texttt{delete\_edge}). Output: a candidate
Python file. Error: a malformed or missing output is a failed call; the loop
continues. & Must & \texttt{test\_prompt\_contains\_pack\_brief},
\texttt{test\_problem\_prompt\_has\_state\_contract},
\texttt{test\_mock\_run\_end\_to\_end} \\
\addlinespace[2pt]

\texttt{FR-AGENT-03} & \textbf{Repair role.} Input: a failed candidate and its
reflection. Precondition: a reflection exists. Behaviour: produce a repaired
candidate. Output: a repaired Python file. Error: as in \texttt{FR-AGENT-02}. &
Must & \texttt{test\_mock\_run\_end\_to\_end}, \texttt{test\_no\_reflection\_ablation} \\
\addlinespace[2pt]

\texttt{FR-AGENT-04} & \textbf{Reflection role.} Input: a failing case (graph,
operation sequence, violated invariant, expected versus actual). Precondition: a
candidate failed. Behaviour: produce a structured reflection with the failure
description, root cause, and repair hints, feeding the next generation's repair.
Output: a reflection file. Error: if reflection fails, the run proceeds without
it (the ablation shows the effect). & Must & \texttt{test\_no\_reflection\_ablation},
\texttt{test\_failing\_context\_survives\_malformed\_state} \\
\addlinespace[2pt]

\texttt{FR-AGENT-05} & \textbf{Prompt and call logging.} Input: every model call.
Precondition: none. Behaviour: log the full prompt, response, tool calls, usage,
and duration for each call, paired with the run's engine file. Output: prompt logs
under the run directory. Error: N/A. & Should &
\texttt{test\_agent\_calls\_paired\_and\_engine\_file\_lifecycle} \\
\addlinespace[2pt]

\texttt{FR-AGENT-06} & \textbf{Agent tool loop.} Input: a role task. Precondition:
the provider is reachable. Behaviour: run a bounded tool loop (read and write the
candidate, run the public examples) up to a step and wall-clock limit, with
reasoning off by default and a nudge when the output is truncated. Output: a
completed or failed call. Error: a step or time limit ends the call as a failed
call, not a failed run. & Must & \texttt{test\_api\_loop\_writes\_checks\_and\_fixes},
\texttt{test\_api\_loop\_step\_limit\_and\_errors},
\texttt{test\_api\_thinking\_off\_and\_length\_nudge} \\
\addlinespace[2pt]

\texttt{FR-AGENT-07} & \textbf{Lean formalization (deferred).} Input: a candidate
that passed the final grade and its problem. Precondition: the experimental flag
is set in a pack. Behaviour: dispatch a Lean prover agent to attempt an honest
proof. Output: a proof attempt that never fails the run. \emph{Status: deferred.}
This step is experimental, requires an external Lean project and harness, and is
not part of the committed contributions; it is stated here so the requirement is
visible and is not presented as met. & Could & \texttt{formalize.py} (present,
not gated); no campaign evidence \\
\bottomrule
\end{longtable}

\subsection{Evolutionary loop}
\label{sec:fr-evo}

\begin{longtable}{@{}p{0.10\textwidth}p{0.40\textwidth}p{0.07\textwidth}p{0.33\textwidth}@{}}
\caption{Functional requirements: evolutionary loop.}\label{tab:fr-evo}\\
\toprule
\textbf{ID} & \textbf{Requirement} & \textbf{Priority} & \textbf{Acceptance evidence} \\
\midrule
\endfirsthead
\toprule
\textbf{ID} & \textbf{Requirement} & \textbf{Priority} & \textbf{Acceptance evidence} \\
\midrule
\endhead
\bottomrule
\endfoot

\texttt{FR-EVO-01} & \textbf{Tournament parent selection.} Input: the scored
population. Precondition: a generation is producing offspring. Behaviour: select
parents by a tournament of size 3 (or the number of viable candidates) over
\texttt{PASS} and \texttt{FAIL} candidates ranked by passing, then fitness, then
lower work ratio. Output: the chosen parents. Error: with no viable candidate,
select none (a fresh generation). & Must &
\texttt{test\_resume\_draws\_the\_same\_selection\_stream},
\texttt{test\_best\_is\_best\_pass\_so\_far} \\
\addlinespace[2pt]

\texttt{FR-EVO-02} & \textbf{Elitism.} Input: the scored population. Precondition:
a generation completed. Behaviour: keep the top-\texttt{elitism\_k} candidates by
rank into the next generation. Output: the surviving elite. Error: a population
smaller than the elitism count is rejected at configuration. & Must &
\texttt{test\_population\_below\_elitism\_is\_rejected} \\
\addlinespace[2pt]

\texttt{FR-EVO-03} & \textbf{Cull with failure deduplication.} Input: the
population above the maximum size. Precondition: a generation completed.
Behaviour: reduce to the maximum size keeping the elite and passing candidates,
and deduplicate failing candidates by failure signature so the population
explores distinct failure modes. Output: the culled population and the discarded
candidates. Error: N/A. & Must & \texttt{test\_best\_is\_best\_pass\_so\_far}
(ranking key); \texttt{population.py} \texttt{cull} \\
\addlinespace[2pt]

\texttt{FR-EVO-04} & \textbf{Stop on success.} Input: the scored population.
Precondition: a generation completed and early stop is enabled. Behaviour: stop
the search when any candidate reaches fitness 1 on the search grade. Output: the
run ends and proceeds to the final grade. Error: with early stop disabled, the
loop continues past success. & Must & \texttt{test\_final\_and\_finish},
\texttt{test\_stop\_before\_start\_ends\_without\_model\_calls} \\
\addlinespace[2pt]

\texttt{FR-EVO-05} & \textbf{Stop on maximum generations.} Input: the generation
counter. Precondition: none. Behaviour: stop when the generation count reaches
\texttt{max\_generations}. Output: the run ends with the best candidate found.
Error: N/A. & Must & \texttt{test\_stop\_mid\_generation\_commits\_what\_was\_produced} \\
\addlinespace[2pt]

\texttt{FR-EVO-06} & \textbf{Stop on consecutive all-failed generations.} Input:
the per-generation outcomes. Precondition: the loop is active. Behaviour: count
consecutive generations in which every candidate failed or crashed; stop when the
count reaches \texttt{max\_consecutive\_all\_failed} (default 5). Output: the run
stops and reports. Error: a single all-failed generation is not fatal. & Should &
\texttt{orchestrator.py} \texttt{\_check\_stop\_conditions};
\texttt{test\_stop\_mid\_generation\_commits\_what\_was\_produced} \\
\addlinespace[2pt]

\texttt{FR-EVO-07} & \textbf{Stop on model-call budget.} Input: the model-call
count. Precondition: none. Behaviour: stop when the count reaches
\texttt{model\_call\_budget} (default 500). Output: the run ends and reports the
budget stop. Error: a crashed generation's calls are still counted on resume. &
Must & \texttt{test\_stop\_before\_start\_ends\_without\_model\_calls},
\texttt{test\_resume\_counts\_calls\_of\_a\_crashed\_generation} \\
\addlinespace[2pt]

\texttt{FR-EVO-08} & \textbf{Plateau detection.} Input: the per-generation best
history. Precondition: dynamic agents are enabled. Behaviour: declare a plateau
when, over the last \texttt{plateau\_window} generations, the population best
improved neither its fitness by \texttt{plateau\_epsilon} nor its work ratio by
\texttt{plateau\_epsilon} (relative). Output: a plateau signal with a signature
(correctness or efficiency). Error: a window shorter than the history does not
trigger. & Must & \texttt{test\_plateau\_on\_fitness\_or\_work\_ratio} \\
\addlinespace[2pt]

\texttt{FR-EVO-09} & \textbf{Spawn a specialist agent.} Input: a plateau signal.
Precondition: fewer than \texttt{agent\_budget} specialists are active and the
time since the last spawn is at least the window. Behaviour: spawn a specialist
generator whose focus is written from the plateau signature; it derives one extra
candidate per generation from the population best. Output: an
\texttt{agent\_spawned} event and the specialist's candidates. Error: with no
plateau, no specialist is spawned. & Must & \texttt{test\_mock\_run\_spawns\_and\_retires\_a\_specialist},
\texttt{test\_spawn\_signature\_and\_retire}, \texttt{test\_focus\_is\_in\_the\_task\_message} \\
\addlinespace[2pt]

\texttt{FR-EVO-10} & \textbf{Retire a specialist agent.} Input: a specialist's
history. Precondition: the specialist is active. Behaviour: retire the specialist
after \texttt{specialist\_generations} candidates or after the window without
producing a new population best. Output: an \texttt{agent\_retired} event with the
reason. Error: a redone generation counts each specialist candidate once. & Must &
\texttt{test\_spawn\_signature\_and\_retire},
\texttt{test\_redone\_generation\_counts\_specialist\_candidates\_once} \\
\bottomrule
\end{longtable}

\subsection{Persistence and recovery}
\label{sec:fr-per}

\begin{longtable}{@{}p{0.10\textwidth}p{0.40\textwidth}p{0.07\textwidth}p{0.33\textwidth}@{}}
\caption{Functional requirements: persistence and recovery.}\label{tab:fr-per}\\
\toprule
\textbf{ID} & \textbf{Requirement} & \textbf{Priority} & \textbf{Acceptance evidence} \\
\midrule
\endfirsthead
\toprule
\textbf{ID} & \textbf{Requirement} & \textbf{Priority} & \textbf{Acceptance evidence} \\
\midrule
\endhead
\bottomrule
\endfoot

\texttt{FR-PER-01} & \textbf{Checkpoint every generation.} Input: a completed
generation. Precondition: none. Behaviour: write a per-generation checkpoint with
the population, the stop reason, the model-call count, and a SHA-256 checksum.
Output: a checkpoint under the run directory. Error: on write failure, mark the
checkpoint incomplete and report it. & Must &
\texttt{test\_stop\_mid\_generation\_commits\_what\_was\_produced},
\texttt{test\_backfill\_existing\_run} \\
\addlinespace[2pt]

\texttt{FR-PER-02} & \textbf{Atomic writes.} Input: any artifact. Precondition:
none. Behaviour: write to a temporary file and rename into place so a crash never
leaves a half-written artifact. Output: a complete artifact or the previous one.
Error: a torn write is impossible by construction. & Must &
\texttt{checkpoint.py} (temp-then-rename); \texttt{test\_malformed\_state\_fails\_cleanly} \\
\addlinespace[2pt]

\texttt{FR-PER-03} & \textbf{Event log for monitoring and replay.} Input: run
events. Precondition: none. Behaviour: append every run event to a line-delimited
JSON log with a monotonic sequence number; the UI and replay read from it.
Output: an \texttt{events.jsonl} that survives restarts. Error: a partial or
malformed line is skipped, not fatal. & Must & \texttt{test\_run\_emits\_full\_event\_stream},
\texttt{test\_seq\_continues\_after\_reopen},
\texttt{test\_partial\_and\_malformed\_lines\_are\_skipped} \\
\addlinespace[2pt]

\texttt{FR-PER-04} & \textbf{Refuse a corrupt checkpoint.} Input: a checkpoint.
Precondition: a resume is requested. Behaviour: validate the checkpoint's
checksum and schema; on mismatch, refuse to resume from it and use the last good
checkpoint. Output: a clean resume or a reported refusal. Error: a corrupt state
is never silently accepted. & Must & \texttt{test\_malformed\_state\_fails\_cleanly},
\texttt{test\_failing\_context\_survives\_malformed\_state} \\
\bottomrule
\end{longtable}

\subsection{Harness (external agent mode)}
\label{sec:fr-har}

\begin{longtable}{@{}p{0.10\textwidth}p{0.40\textwidth}p{0.07\textwidth}p{0.33\textwidth}@{}}
\caption{Functional requirements: harness.}\label{tab:fr-har}\\
\toprule
\textbf{ID} & \textbf{Requirement} & \textbf{Priority} & \textbf{Acceptance evidence} \\
\midrule
\endfirsthead
\toprule
\textbf{ID} & \textbf{Requirement} & \textbf{Priority} & \textbf{Acceptance evidence} \\
\midrule
\endhead
\bottomrule
\endfoot

\texttt{FR-HAR-01} & \textbf{Initialize a workspace.} Input: a problem pack and a
target directory. Precondition: the directory is empty. Behaviour: create a run
folder and the agent files (\texttt{workspace.json}, \texttt{program.md},
\texttt{candidate.py}, \texttt{results.tsv}). Output: a workspace an external
coding agent can work in. Error: a name/folder mismatch, a missing brief, or a
non-empty directory is rejected. & Must & \texttt{test\_init\_creates\_agent\_files\_and\_run},
\texttt{test\_name\_folder\_mismatch\_rejected}, \texttt{test\_missing\_brief\_rejected},
\texttt{test\_check\_finds\_workspace\_from\_subdir\_and\_rejects\_bad\_level} \\
\addlinespace[2pt]

\texttt{FR-HAR-02} & \textbf{Grade the workspace candidate.} Input: a check
request. Precondition: a workspace exists. Behaviour: grade only the workspace's
\texttt{candidate.py}, store a hashed snapshot, append a row to
\texttt{results.tsv}, and emit the usual events so the UI shows the session like a
run. Output: a grade (examples, search, or full level). Error: a file other than
\texttt{candidate.py} is not graded. & Must & \texttt{test\_check\_grades\_only\_the\_workspace\_candidate},
\texttt{test\_check\_records\_like\_a\_run} \\
\addlinespace[2pt]

\texttt{FR-HAR-03} & \textbf{Independent final grade.} Input: a kept candidate.
Precondition: the agent requests \texttt{check --final}. Behaviour: run the
independent final grade on the full bound; this is what counts for evaluation.
Output: a final-grade verdict. Error: a final grade is allowed even after
\texttt{max\_checks} is reached, but only once. & Must &
\texttt{test\_check\_accepts\_candidate\_by\_other\_path\_and\_one\_final},
\texttt{test\_max\_checks\_is\_enforced\_but\_final\_still\_runs} \\
\addlinespace[2pt]

\texttt{FR-HAR-04} & \textbf{Enforce harness integrity.} Input: the workspace
state. Precondition: none. Behaviour: enforce \texttt{max\_checks}, detect later
edits to a graded snapshot by its hash, refuse further checks on a finished
workspace, and allow-list the harness environment. Output: an integrity-enforced
session. Error: a snapshot edit, an over-limit check, or a check on a finished
workspace is refused. & Must & \texttt{test\_graded\_snapshot\_hash\_detects\_later\_edits},
\texttt{test\_finished\_workspace\_refuses\_more},
\texttt{test\_harness\_env\_is\_allow\_listed} \\
\addlinespace[2pt]

\texttt{FR-HAR-05} & \textbf{Finish a workspace run.} Input: a finish request.
Precondition: a workspace exists. Behaviour: close the run, emitting a final
\texttt{run\_end}; an idle harness run counts as gone. Output: a closed run the UI
shows as finished. Error: finishing an already-finished workspace is a no-op. &
Should & \texttt{test\_final\_and\_finish}, \texttt{test\_idle\_harness\_run\_counts\_as\_gone} \\
\bottomrule
\end{longtable}
